\section{Cross-sample prediction fragility}
\label{sec:preliminaries}

\paragraph{Setting.}
A learner $\mathcal{A}$ maps a training set
$S = \{(x_i, y_i)\}_{i=1}^{N}$ drawn iid from population $\mathcal{D}$
to a classifier $f_S = \mathcal{A}(S)$.  At deployment time the
practitioner reports $\hat{y}(x) = \arg\max_c f_S(x)_c$ for each test
example $x$.  We are interested in how much $\hat{y}(x)$ changes when
$S$ is replaced by an independent draw from $\mathcal{D}$.

\paragraph{Definition.}
Given two iid training samples $S_A, S_B \sim \mathcal{D}^N$ and a
test example $x$, the \emph{cross-sample fragility} of $\mathcal{A}$ at
$x$ is
\begin{equation}
  \rho(\mathcal{A}, x) =
  \mathbb{P}_{S_A, S_B}\!\big[\arg\max f_{S_A}(x) \neq \arg\max f_{S_B}(x)\big].
\end{equation}
A fragility-aware report supplements the standard accuracy with the
expected per-example argmax disagreement
$\bar\rho(\mathcal{A}) = \mathbb{E}_x [\rho(\mathcal{A}, x)]$, and
optionally with a distributional analogue
$\bar\rho_{\mathrm{KL}}(\mathcal{A})
= \mathbb{E}_x \big[\tfrac{1}{2}(\mathrm{KL}(f_{S_A}(x) \| f_{S_B}(x))
                                + \mathrm{KL}(f_{S_B}(x) \| f_{S_A}(x)))\big]$.
Throughout the paper, ``fragility'' refers to $\bar\rho$ unless
otherwise stated, and ``argmax churn'' is the same quantity reported
as a percentage; figures use ``id-churn'' on axes when the canonical
id-test set is the population the expectation is taken over.

\paragraph{Bootstrap estimator.}
In practice we have one observed training set $S$, so we estimate
$\rho$ by drawing two independent bootstraps from $S$ (sample $N$
items with replacement).  Each bootstrap is a draw of size $N$ from
the empirical distribution $\hat{\mathcal{D}}_S$, and $\hat{\mathcal{D}}_S
\to \mathcal{D}$ as $N \to \infty$.  We treat the bootstrap-pair
disagreement as a proxy for $\rho$ throughout this paper; we do not
claim a formal consistency result.  All numerical magnitudes we report
are with respect to this bootstrap proxy, which is what a practitioner
can compute from their own training set.

\paragraph{Connection to algorithmic stability.}
\citet{TODO_bousquet2002} define uniform stability $\beta$ on the
loss when a single training example is altered.  Cross-sample
fragility is the corresponding output-space quantity under
multi-example perturbation, and decreases with $N$ for the same
reason $\beta$ does.  We use $\beta$ as the interpretive anchor but
do not claim fragility matches the $O(1/N)$ rate quantitatively;
Appendix~\ref{app:theory} enumerates the gaps.

\paragraph{Strict measurement protocol.}
To avoid contaminating cross-sample measurements with test-set drift, we
fix a canonical test set once and only vary the training-side seed.
Concretely: a single canonical seed $99$ constructs the test set; for
each measurement seed $s \in \{1, \dots, 10\}$ we draw a bootstrap
$S^{(s)}$ from the canonical training set and train
$f^{(s)} = \mathcal{A}(S^{(s)})$.  Cross-sample churn is the average of
$\mathbf{1}[\arg\max f^{(s)}(x) \neq \arg\max f^{(s')}(x)]$ over all
$\binom{10}{2}$ seed pairs.  Confidence intervals are obtained by
bootstrap resampling of seed pairs; deltas between methods are paired.
